\documentclass[11pt,a4paper]{article}
\usepackage[T1]{fontenc}
\usepackage{lmodern}
\usepackage{amsmath,amssymb,amsfonts}
\usepackage{geometry}
\geometry{a4paper, margin=25mm}
\usepackage{booktabs}
\usepackage{xcolor}
\usepackage{fancyhdr}
\usepackage{hyperref}
\usepackage[style=numeric,sorting=none]{biblatex}
\addbibresource{refs.bib}

\hypersetup{
    colorlinks=true,
    linkcolor=blue!60!black,
    citecolor=green!50!black,
    urlcolor=blue!80!black
}

\pagestyle{fancy}
\fancyhf{}
\rhead{\small CS/DT Coursework 2026}
\lhead{\small Deliverable 1: High-Level Programmatic Document}
\cfoot{\thepage}

\title{\textbf{Project Specification: Layer-Resolved Structural Aliasing\\
and Cyber-Physical Vulnerabilities in Time-Series Foundation Models}}
\author{\textbf{Deliverable 1: High-Level Programmatic Document}\\
Joint Computer Security / Digital Transformation Graduate Program\\
Coursework Co-Supervision: Prof. Benini\\
Project Collaborators: Brignoli, Boniotti, and Sabbadini}
\date{\today}

\begin{document}
\maketitle

%-----------------------------------------------------------------------
\section{Purpose and Scope (The ``What?'')}
%-----------------------------------------------------------------------

This project formalises and empirically validates the phenomenon of \textit{structural aliasing} in patch-based time-series foundation models, using Amazon Chronos-Bolt-Tiny as the target architecture~\cite{ansari2024chronos}. Chronos-Bolt-Tiny operates with a fixed patch size of $P=16$ and a default stride of $S=16$, a context window of $T=512$ tokens, and four encoder blocks followed by four decoder blocks terminating in a direct quantile projection head. The baseline evaluation uses $S=16$; stride is then systematically varied as the independent variable to test hypothesis H3 (Section~3).

Layer-resolved probing is performed exclusively on the four encoder hidden states (enc\,0--3), where each stage produces a sequence of tokens amenable to representational analysis. The decoder, which operates on a single query token, is analysed only at the level of its output probability distribution. This scope is therefore explicitly bounded to \emph{encoder-stage representational analysis plus output-level behavioural evaluation}; it excludes model retraining, fine-tuning, and comparative evaluation of alternative foundation model architectures.

The empirical anchor is the 1-minute multivariate telemetry stream in \texttt{Dataset-SolarTechLab.csv}, prioritising active photovoltaic power generation ($P_{\mathrm{PV}}$, target column \texttt{Power\_W}) and global tilted irradiance ($G_{\mathrm{tilt}}$) where available. This domain provides naturally recurrent signals whose dominant periods (diurnal, cloud-intermittency) span both safe and architecturally blind frequency bands, making structural aliasing directly observable.

%-----------------------------------------------------------------------
\section{Intended Audience and Impact (The ``Why?'')}
%-----------------------------------------------------------------------

The primary audience is the academic evaluation committee led by Prof. Benini, together with researchers working on time-series representation learning and interpretability. A secondary audience comprises smart-grid operations engineers who deploy pre-trained forecasting models in cyber-physical control loops.

The core objective is to challenge the dangerous assumption that Nyquist-compliant signals are safely processed by patch-based temporal foundation models. We demonstrate that when a signal's period is commensurate with the model's patch geometry, an \emph{architectural data-deletion} event renders the signal invisible to self-attention, irrespective of sampling-theorem compliance~\cite{oppenheim2010}. In automated micro-grid settings, this representation collapse propagates to semantic misclassifications, with concrete downstream consequences including battery over-discharge, flawed natural-language anomaly explanations, and unmitigated cyber-physical instability. Framing this mechanism as a formal vulnerability within the NIST AI Risk Management Framework~\cite{nist2023} provides practitioners with actionable risk-management language, and directly addresses the security requirements of the joint CS/DT track.

%-----------------------------------------------------------------------
\section{Persuasion Strategy (The ``How?'')}
%-----------------------------------------------------------------------

\paragraph{1. Formal derivation.}
We model patch-based tokenisation as a linear operator $T_{P,S}:\mathbb{R}^N\!\to\mathbb{R}^{\lceil(N-P)/S+1\rceil\times P}$ and derive analytically when two consecutive patches become identical. The key quantity is cycles per patch $\mathrm{cpp}=fP/f_s$: phase-locking with the patch grid occurs at integer cpp values ($\mathrm{cpp}\in\mathbb{Z}^+$), and the token-Nyquist ceiling is $\mathrm{cpp}=P/2=8$. These closed-form conditions constitute the theoretical backbone against which all empirical results are interpreted.

\paragraph{2. Falsifiable hypotheses.}
Three hypotheses are formally stated and independently tested:
\begin{itemize}
    \item \textbf{H1} (Blind spots): systematic representational collapse occurs at frequencies satisfying $\mathrm{cpp}\in\mathbb{Z}^+$, and not at adjacent off-grid frequencies.
    \item \textbf{H2} (Phase sensitivity): at any blind-spot frequency, patch-sequence redundancy is phase-invariant---the effect is independent of the signal's temporal offset relative to the patch grid.
    \item \textbf{H3} (Geometry dependence): reducing stride $S < P$ shifts blind-spot frequencies and reduces redundancy; increasing $P$ lowers the blind-spot band.
\end{itemize}

\paragraph{3. Dual-engine synthetic signal generation.}
Controlled aliasing experiments use \textit{KernelSynth}~\cite{ansari2024chronos} (composite Gaussian-process prior, sampling-rate dependent) as the primary instrument: deterministic sinusoids are injected at known cpp values into a stochastic background, isolating architectural effects from classical sampling artefacts. \textit{TSMixup} (Dirichlet mixup of real PV subsequences, sampling-rate agnostic) serves exclusively as a confound control, verifying that observed representational effects are not artefacts of the synthetic generation process.

\paragraph{4. Layer-resolved probing on encoder states.}
We deploy Minimum Description Length (MDL) probes~\cite{voita2020} on the hidden states extracted after each of the four encoder blocks, quantifying the linear decodability of frequency attributes via the MDL compression ratio. Decoder analysis is restricted to output-distribution statistics (quantile spread, spectral RMSE between forecast and ground truth). Causal intervention applies LEACE~\cite{belrose2023} to surgically erase coarse frequency abstractions from encoder representations and measures the downstream impact on forecasting accuracy.

\paragraph{5. Semantic ontology layer.}
Domain concepts (e.g., \textit{diurnal\_cycle}, \textit{cloud\_intermittency}, \textit{night\_zero}) are formalised in a PV ontology mapping each concept to its characteristic cpp value and aliasing-risk category. This layer makes the link between signal structure, learned representation, and domain semantics explicit, and identifies which operational states are observationally invisible to the model at native 1-minute sampling.

\paragraph{6. Bayesian statistical analysis.}
All effect estimates are obtained through a probabilistic model implemented in PyMC. A \textit{Skeptical Prior} $\beta\sim\mathcal{N}(0,\,0.1^2)$ encodes prior scepticism about the aliasing effect size; an observation-noise prior $\sigma\sim\mathrm{HalfNormal}(0.0616)$ is calibrated from the empirical standard deviation of \texttt{Power\_W} during nighttime null intervals in the dataset, representing baseline instrument noise under zero-generation conditions. Posterior distributions, 95\% credible intervals, and Bayes factors comparing the aliasing model to a null (frequency-flat) model are reported throughout.

\paragraph{7. Mitigation.}
We evaluate randomised patch-offset jitter as the primary mitigation: adding a uniform random offset $\delta\sim\mathcal{U}(0,P-1)$ to the patch start index breaks phase-locking without architectural changes. The Bayesian framework is reused to quantify residual risk post-mitigation.

\paragraph{8. Agentic AI and NIST AI RMF.}
We discuss how a monitoring agent could detect, explain, and act upon aliasing risks in real time---for instance, by maintaining a cpp map of live signals and flagging tokens with high blind-band energy exposure. The identified data-deletion mechanism is framed as a systematic architectural vulnerability within all four core functions of the NIST AI RMF 1.0~\cite{nist2023} (\textit{Govern, Map, Measure, Manage}), evaluating whether the proposed mitigation reduces residual risk or leaves significant exposure.

%-----------------------------------------------------------------------
\section{Open Source and Licensing Compliance}
%-----------------------------------------------------------------------

This document and the final report are released under CC-BY 4.0. All source code, probing scripts, and intervention wrappers are released under the MIT licence. Any third-party material is used in compliance with its own licence and clearly attributed.

%-----------------------------------------------------------------------
\printbibliography
%-----------------------------------------------------------------------
\end{document}
